\subsection{通常表示簇的实结构}
\label{sec:02.3.2}

为了能对上面的映射 $\Phi$ 说些什么，我们需要固定一个额外的结构，即 $V$ 上的线性实
结构 $\alpha$。特别地，由引理 2.3.1，群 $\U^{\pm}(V)$ 现在典范同构于半直积
$\U(V)\rtimes_\alpha\Z/2\Z$。我们研究的第一步，是弄清 $\pi_1(X;x)$ 的酉表示簇上的
实结构是什么。然后，作为第二步，我们将证明 $\Phi$ 把
$\pi^R_1((X,\sigma);x)$ 的表示簇送到 $\pi_1(X;x)$ 的表示簇的实点。

短正合列
\[
1\longrightarrow\pi_1(X;x)\longrightarrow\pi^R_1((X,\sigma);x)
\longrightarrow\Z/2\Z\longrightarrow1
\]
的存在性蕴涵群同态
\[
\Z/2\Z\longrightarrow\Out(\pi_1(X;x)):=
\Aut(\pi_1(X;x))/\Inn(\pi_1(X;x))
\]
的存在性。事实上，设 $f\in\pi^R_1((X,\sigma);x)$ 在增广映射下送到
$\sigma\in\Z/2\Z$。则 $f$ 通过共轭作用在 $\pi^R_1((X,\sigma);x)$ 的正规子群
$\pi_1(X;x)$ 上。纤维 $\sigma$ 中另一元素 $f'$ 的任一其他选取与 $f$ 相差
$\pi_1(X;x)$ 中元素 $f^{-1}f'$，故与 $f$ 诱导的 $\pi_1(X;x)$ 自同构相差一个内自
同构（事实上，$\Ad_f^{-1}\circ\Ad_{f'}=\Ad_{f^{-1}f'}$，其中
$f^{-1}f'\in\pi_1(X;x)$），于是我们有映射
$\Z/2\Z\to\Out(\pi_1(X;x))$。这个映射确实是群同态（关于此映射的更多信息见习题
2.3.7 与 2.3.9）。现在回忆：群 $\Out(\pi_1(X;x))$ 通过
\begin{equation}\label{eq:2.23}
[\phi]\cdot[\rho]=[\rho\circ\phi^{-1}]
\end{equation}
左作用于表示簇 $\Hom(\pi_1(X;x);\U(V))/\U(V)$，其中 $[\phi]$ 是
$\phi\in\Aut(\pi_1(X;x))$ 的外类；见习题 2.3.8。这个作用称为\emph{外作用}。在我们的
实代数情境中，由于有映射
$\Z/2\Z\to\Out(\pi^R_1((X,\sigma);x))$，它给出 $\Z/2\Z$ 在 $\pi_1(X;x)$ 的表示簇上
的作用。确切地说，若 $\tilde\sigma$ 位于增广映射
$\pi^R_1((X,\sigma);x)\to\Z/2\Z$ 的纤维 $\sigma$ 之上，则相应的在
$\Hom(\pi_1(X;x);\U(V))/\U(V)$ 上的作用由
$[\rho]\mapsto[\rho\circ\Ad_{\tilde\sigma}^{-1}]$ 定义，其中
$\Ad_{\tilde\sigma}^{-1}=\Ad_{\tilde\sigma^{-1}}\in\Aut(\pi_1(X;x))$ 是在
$\pi_1(X;x)$ 中通过 $\tilde\sigma^{-1}$ 的共轭（我们回忆：$\tilde\sigma$ 在
$\pi^R_1((X,\sigma);x)$ 中一般不是二阶元，故上式中的 $-1$ 指数要紧）。所得表示簇中
的元素 $[\rho\circ\Ad_{\tilde\sigma}^{-1}]$ 不依赖于纤维 $\sigma$ 中 $\tilde\sigma$ 的
选取；又因 $\tilde\sigma^2\in\pi_1(X;x)$，我们得到 $\Z/2\Z$ 在
$\Hom(\pi_1(X;x);\U(V))/\U(V)$ 上的作用。$\pi_1(X;x)$ 的酉表示簇的这个对合以纯代数
方式定义，没有用到目标群 $\U(V)$ 的任何额外结构。现在回忆：后者具有来自 $V$ 上实结构
$\alpha$ 的 Galois 共轭
\[
\Ad_\alpha:\ \U(V)\longrightarrow\U(V),\qquad u\longmapsto\alpha\circ u\circ\alpha^{-1}
\]
（复合在 $\U^{\pm}(V)$ 中进行，但所得元素确实落在 $\U(V)$ 中）。于是我们有对合
\[
\Hom(\pi_1(X;x);\U(V))/\U(V)\longrightarrow\Hom(\pi_1(X;x);\U(V))/\U(V),
\qquad [\rho]\longmapsto[\Ad_\alpha\circ\rho].
\]
通常记 $\bar\rho:=\Ad_\alpha\circ\rho$（我们已看到，在 $V$ 的适当基下用实结构 $\alpha$
共轭矩阵所做的正是复共轭）。注意 $\rho\mapsto\bar\rho$ 的不动点特别地保持实子空间
$V^\alpha\subset V$：换言之，当 $\rho=\bar\rho$ 时表示 $\rho$ 取值于
$\GL(V^\alpha)$（不过一般来说，$\Hom(\pi_1(X;x);\U(V))/\U(V)$ 的诱导对合的不动点并非
都来自这样的 $\rho$）。现在我们把刚才考虑的两个对合合起来，得到表示簇的一个实结构，
$\pi^R_1((X,\sigma);x)$ 的表示关于它的表现十分特别（它们是这个实结构的实点）。注意
我们把它称为实结构，是因为可以证明它关于表示簇的 Narasimhan--Seshadri 复结构是反
全纯的（例如可参看 [16]）；但在现阶段，它只是一个对合。我们用
\[
\mathcal{R}:=\Hom(\pi_1(X;x);\U(V))/\U(V)
\]
记 $\pi_1(X;x)$ 的酉表示簇，用
\begin{equation}\label{eq:2.24}
\mathcal{R}(\sigma,\alpha):=
\Hom_{\Z/2\Z}\bigl(\pi^R_1((X,\sigma);x);\U(V)\rtimes_\alpha\Z/2\Z\bigr)/\U(V)
\end{equation}
记 $\pi^R_1((X,\sigma);x)$ 的酉表示簇，并用
\begin{equation}\label{eq:2.25}
\kappa:\ \mathcal{R}\longrightarrow\mathcal{R},
\qquad [\rho]\longmapsto\bigl[\Ad_\alpha\circ\rho\circ\Ad_{\tilde\sigma}^{-1}\bigr]
\end{equation}
记由前述两个对合复合得到的表示簇 $\mathcal{R}$ 的对合。注意 $\kappa$ 同时依赖于
$\sigma$ 与 $\alpha$。

\begin{theorem}\label{thm:2.3.4}
映射 $\Phi:\mathcal{R}(\sigma,\alpha)\to\mathcal{R}$，
$[\rho]\mapsto\bigl[\rho|_{\pi_1(X;x)}\bigr]$ 把 $\mathcal{R}(\sigma,\alpha)$ 送到
$\Fix(\kappa)\subset\mathcal{R}$ 的点；一般地它不是满射。
\end{theorem}

\begin{proof}
对合 $\kappa$ 把群同态 $\chi:\pi_1(X;x)\to\U(V)$ 的类送到
\[
\Ad_\alpha\circ\chi\circ\Ad_{\tilde\sigma}^{-1}:\ \pi_1(X;x)\longrightarrow\U(V),
\qquad f\longmapsto\alpha\chi(\tilde\sigma^{-1}f\tilde\sigma)\alpha^{-1}
\]
的类。但若 $\chi$ 落在 $\Phi$ 的像中，则 $\chi$ 可延拓为群同态
$\rho:\pi^R_1((X,\sigma);x)\to\U(V)\rtimes_\alpha\Z/2\Z$。于是
\[
\begin{aligned}
\Ad_\alpha\circ\chi\circ\Ad_{\tilde\sigma}^{-1}(f)
&=\Ad_\alpha\circ\rho\circ\Ad_{\tilde\sigma}^{-1}(f)\\
&=\alpha\bigl(\rho(\tilde\sigma^{-1})\rho(f)\rho(\tilde\sigma)\bigr)\alpha^{-1}\\
&=\bigl(\alpha\rho(\tilde\sigma)^{-1}\bigr)\rho(f)
\bigl(\alpha\rho(\tilde\sigma)^{-1}\bigr)^{-1}\\
&=\underbrace{\bigl(\alpha\rho(\tilde\sigma)^{-1}\bigr)}_{\in\,\U(V)}
\chi(f)
\underbrace{\bigl(\alpha\rho(\tilde\sigma)^{-1}\bigr)^{-1}}_{\in\,\U(V)}.
\end{aligned}
\]
事实上，$\alpha$ 与 $\rho(\tilde\sigma)$ 都落在 $\U(V)\rtimes_\alpha\Z/2\Z$ 中纤维
$\sigma$ 之上，故 $\alpha\rho(\tilde\sigma)^{-1}$ 落在平凡元之上的纤维中，从而属于
$\U(V)$。因此 $\bigl(\Ad_\alpha\circ\chi\circ\Ad_{\tilde\sigma}^{-1}\bigr)
\sim_{\U(V)}\chi$，这意味着 $[\chi]\in\Fix(\kappa)$。关于一般地
$\mathrm{Im}\,\Phi\ne\Fix(\kappa)$ 的证明，可参看 [16]。
\end{proof}

\subsection*{习题}

\begin{exercise}\label{ex:2.3.5}
证明：复向量空间 $V$ 上（线性）实结构之集与齐性空间
$\GL(n;\C)/\GL(n;\R)$ 双射，其中 $n$ 是 $V$ 的复维数。
\end{exercise}

\begin{exercise}\label{ex:2.3.6}
证明引理 2.3.1。并证明：若 $\C$-线性变换 $u$ 在 $(V,\alpha)$ 的某个给定实基下的
矩阵为 $A$，则 $\alpha\circ u\circ\alpha^{-1}$ 的矩阵为 $\bar A$，即 $A$ 的复共轭。
\end{exercise}

\begin{exercise}\label{ex:2.3.7}
设 $G$ 是群，$g\in G$，$f\in\Aut(G)$。证明 $f\circ\Ad_g\circ f^{-1}=\Ad_{f(g)}$，
并证明每个短正合列 $1\to G\to H\to I\to1$ 都诱导群同态
$I\to\Out(G):=\Aut(G)/\Inn(G)$，其中
$\Inn(G):=\{\Ad_g\mid g\in G\}$ 是 $G$ 的内自同构群。
\end{exercise}

\begin{exercise}\label{ex:2.3.8}
验证外作用 (2.23) 是良定义的。
\end{exercise}

\begin{exercise}\label{ex:2.3.9}
证明：若 $\sigma(x)=x$，则
\begin{enumerate}[label=(\roman*), leftmargin=2.2em, itemsep=0.15em]
\item 同态 $\Z/2\Z\to\Out(\pi_1(X;x))$ 把 $\sigma\in\Z/2\Z$ 送到由
      $\sigma:X\to X$ 诱导的群自同构 $\sigma_*:\pi_1(X;x)\to\pi_1(X;x)$；
\item $\Z/2\Z$ 事实上已经通过 $\rho\mapsto\rho\circ\sigma_*^{-1}$ 左作用于
      $\Hom(\pi_1(X;x);\U(V))$。
\end{enumerate}
\end{exercise}

\begin{exercise}\label{ex:2.3.10}
设 $p:B\to C$ 是满同态，并用
\begin{equation}\label{eq:2.26}
1\longrightarrow A\longrightarrow B\longrightarrow C\longrightarrow1
\end{equation}
记相应的群的短正合列。
\begin{enumerate}[label=(\roman*), leftmargin=2.2em, itemsep=0.15em]
\item 设 $s:C\to B$ 是满足 $p\circ s=\Id_C$ 的群同态（这样的映射称为短正合列 (2.26)
      的\emph{分裂}），考虑相应的半直积 $A\rtimes_s C$，其中
      $(a,c)(a',c')=(as(c)a's(c)^{-1},cc')$。证明映射
\[
\phi_s:\ A\rtimes_s C\longrightarrow B,
\qquad (a,c)\longmapsto a\,s(c)
\]
诱导 $C$ 被 $A$ 的扩张的同构：
\[
\begin{tikzcd}[column sep=large, row sep=large]
A \arrow[r] \arrow[d, "="'] & A\rtimes_s C \arrow[r] \arrow[d, "\phi_s"', "\simeq"] & C \arrow[d, "="]\\
A \arrow[r] & B \arrow[r] & C
\end{tikzcd}
\]
\item 设 $u\in A$。证明 $s':=\Ad_u\circ s$ 是 (2.26) 的另一个分裂，且作为 $C$ 的增广
      有 $A\rtimes_{s'}C\simeq A\rtimes_s C$：
\[
\begin{tikzcd}[column sep=large, row sep=large]
A \arrow[r] \arrow[d, "\Ad_u"'] & B \arrow[r] \arrow[d, "\Ad_u"'] & C \arrow[d, "="]\\
A \arrow[r] & B \arrow[r] & C
\end{tikzcd}
\]
\end{enumerate}
\end{exercise}
